\documentclass[12pt]{report}
\begin{document}
	\sffamily
	\section{Herramientas}
		\subsection{IDEs}
			% \subsubsection{Concepto}
			\subsubsection{Recomendaciones}
				% table con nombre del IDE con el enlace de descarga hiperreferenciado, 
				% breve descripción, ventajas, desventajas, y tal vez imagen
				% \begin{landscape}
				\begin{table}[!hbt]
					\caption{Recomendaciones de IDEs}
					\begin{center}
						\begin{tabular}{c c c}
							\textbf{Nombre}	&	\textbf{Descripción}	&	\textbf{Es el indicado si...}\\
							\hline
							% Atom
								\parbox[c]{2cm}{
									\centering
									\href{https://atom.io}{Atom}
									\newline\newline
									\includegraphics[height=1.5cm]{logo_atom}	
								}	
								&
								\parbox[c]{3cm}{
									Editor de código multiplataforma creado por GitHub
								}
								&	
								\parbox[c]{8cm}{
									\begin{itemize}
										\item Necesitas algo ligero
										\item Quieres algo sencillo de usar y aprender
										\item Quieres trabajar colaborativamente
										\item Quieres personalizar tu IDE
									\end{itemize}
								}
								\\\hline
							% VSC
								\parbox[c]{2cm}{
									\centering
									\href{https://code.visualstudio.com/download}{VS Code}
									\newline\newline
									\includegraphics[height=1.5cm]{logo_vscode}	
								}
								&
								\parbox[c]{3cm}{
									IDE multiplataforma creado por Microsoft
								}
								&
								\parbox[c]{8cm}{
									\begin{itemize}
										\item Vas a trabajar en proyectos grandes
										\item Quieres trabajar con varios lenguajes
										\item Quieres personalizar tu IDE
										\item Quieres usar varias extensiones
									\end{itemize}
								}
								\\\hline
							% JetBrains
								\parbox[c]{2cm}{
									\centering
									\href{https://www.jetbrains.com}{JetBrains}
									\newline\newline
									\includegraphics[height=1.5cm]{logo_jetbrains}	
								}	
								&	
								\parbox[c]{3cm}{
									Es una familia de IDEs especializados con versiones libres y de pago como 
									\href{https://www.jetbrains.com/es-es/phpstorm/}{PHPstorm} o 
									\href{https://www.jetbrains.com/es-es/webstorm/}{WebStorm}
								}
								&
								\parbox[c]{8cm}{
									\begin{itemize}
										\item Vas a trabajar en un área especializada, e.g.: puro backend o puro frontend
										\item Tienes una máquina con buen rendimiento
										\item Quieres ver tus errores en tiempo real % corregir
									\end{itemize}
								}
								\\\hline
							% Repl.it
								\parbox[c]{2cm}{
									\centering
									\href{https://replit.com}{Repl.it}
									\newline\newline
									\includegraphics[height=1.5cm]{logo_replit}	
								}	
								&
								\parbox[c]{3cm}{
									IDE en línea colaborativo inspirado en Google Docs
								}
								&
								\parbox[c]{8cm}{
									\begin{itemize}
										\item No puedes instalar programas
										\item Requieres colaboración eficiente en tiempo real
										\item Quieres todas las ventajas de un IDE sin necesidad de configurarlo
									\end{itemize}
								}
								\\\hline
						\end{tabular}
					\end{center}
				\end{table}
				% \end{landscape}
	% 		\subsubsection{Instalación de XAMPP y MAMP}
	% 			\paragraph{XAMPP}
	% 			\paragraph{MAMP}
	% 	\subsection{Prototipado}
	% 		\subsubsection{Figma}
	% 			\paragraph{¿Qué es y para qué sirve?}
	% 			\paragraph{Menú}
	% 				\begin{enumerate}
	% 					\item Herramientas
	% 					\item ComparacionesPropiedades
	% 				\end{enumerate}
	% 			\paragraph{Herramientas básicas}
	% 				\begin{enumerate}
	% 					\item Frames
	% 					\item FigurasTexto
	% 					\item Pluma
	% 				\end{enumerate}
	% 			\paragraph{Propiedades y capas}
	% 				\begin{enumerate}
	% 					\item Orden
	% 					\item Color
	% 					\item Dimensiones
	% 				\end{enumerate}
	% 			\paragraph{Compartir}
	% 		% tal vez cambiar por una tabla con nombre, descripción, ventajas y desventajas
	% 		\subsubsection{Adobe XD}
	% 		\subsubsection{Sketch}
	% 		\subsubsection{Justinmind}
	% 	\subsection{Controladores de versiones}
	% 		\subsubsection{Git}
	% 			\paragraph{Concepto}
	% 			\paragraph{Creación de repositorios}
	% 			\paragraph{Comandos}
	% 				% table con el nombre (git + clone, status, add, commit, push), descripción, 
	% 				% código con las opciones más importantes, caso de uso, y resultado
	% 		\subsubsection{Páginas para el control de repositorios creados con Git}
	% 			\paragraph{GitHub}
	% 				\subparagraph{Uso de GitHub}
	% 				\subparagraph{Visualización y modificación de repositorios}
	% 				\subparagraph{Integración de git para el manejo de proyectos alojados en GitHub}
	% \section{Seguridad web}
	% 	\subsection{Tipos de ataques informáticos y cómo protegerse de ellos}
	% 		\subsubsection{XSS}
	% 		\subsubsection{CSRF}
	% 		\subsubsection{Inyección SQL}
	% 	\subsection{Validaciones}
	% 		\subsubsection{Contraseñas seguras}
	% 		\subsubsection{Expresiones regulares}
	% 			\paragraph{Uso en validación de datos}
	% 		\subsubsection{Lado del cliente (HTML,JS)}
	% 		\subsubsection{Lado del servidor (PHP)}
	% 	\subsection{Información y datos sensibles}
	% 		\subsubsection{Cifrado}
	% 			\paragraph{Definición}
	% 			\paragraph{Tipos}
	% 				\begin{enumerate}
	% 					\item Cifrado simétrico
	% 					\item Cifrado asimétrico
	% 					\item Cifrado híbrido
	% 				\end{enumerate}
	% 		\subsubsection{Hasheo}
	% 			\begin{enumerate}
	% 				\item Definición
	% 				\item Sal
	% 				\item Pimienta
	% 			\end{enumerate}
	% \section{Arquitectura de la información}
	% 	\subsection{Definición}
	% 	\subsection{Los 3 círculos}
	% 		\subsubsection{Contexto}
	% 		\subsubsection{Contenido}
	% 		\subsubsection{Usuarios}
	% 	\subsection{Encontrando información}
	% 		\subsubsection{Necesidades}
	% 		\subsubsection{Modelos de búsqueda}
	% 		\subsubsection{Retroalimentación}
	% 	\subsection{Entendiendo información}
	% 		\subsubsection{Tipologías}
	% 		\subsubsection{\textit{Capas} del entorno de información}
	% 		\subsubsection{Flexibilidad}
	% 	\subsection{Top-Down AI}
	% 	\subsection{Bottom-Up AI}
	% 		\subsubsection{Stress-test}
	% 	\subsection{Sistemas de Organización}
	% 		\subsubsection{Sistemas de organización}
	% 			\paragraph{Dificultades}
	% 			\paragraph{Esquemas de organización}
	% 				\subparagraph{Exactos}
	% 				\subparagraph{Antiguos}
	% 					\begin{enumerate}
	% 						\item Tópicos
	% 						\item Orientado a tareas
	% 						\item Específico a audiencia
	% 						\item metafóricos
	% 						\item Híbridos
	% 					\end{enumerate}
	% 			\paragraph{Estructuras de organización}
	% 				\begin{enumerate}
	% 					\item Jerarquía
	% 						\begin{enumerate}
	% 							\item Diseñando jerarquías
	% 							\item Amplitud vs. profundidad
	% 						\end{enumerate}
	% 					\item Modelo Relacional
	% 					\item Hipertexto
	% 					\item Clasificación social
	% 				\end{enumerate}
	% 		\subsubsection{Etiquetado de sistemas}
	% 			\paragraph{Importancia}
	% 			\paragraph{Variedad}
	% 				\begin{enumerate}
	% 					\item Enlaces
	% 					\item Encabezados
	% 					\item Opciones de Navegación
	% 					\item Index
	% 					\item Íconos
	% 				\end{enumerate}
	% 			\paragraph{Diseño de etiquetas}
	% 				\begin{enumerate}
	% 					\item Alcance
	% 					\item Consistencia
	% 					\item Pruebas
	% 				\end{enumerate}
	% 		\subsubsection{Sistemas de navegación}
	% 			\paragraph{Flexibilidad}
	% 			\paragraph{Tipos}
	% 				\begin{enumerate}
	% 					\item Locales
	% 					\item Globales
	% 					\item Contextuales
	% 				\end{enumerate}
	% 			\paragraph{Suplementarios}
	% 			\paragraph{Personalización}
	% 			\paragraph{Visualización}
	% 			\paragraph{Social}
	% 		\subsubsection{Sistemas de búsqueda}
	% 			\paragraph{Necesidad}
	% 			\paragraph{Anatomía}
	% 			\paragraph{Filtros}
	% 			\paragraph{Algoritmos}
	% 				\subparagraph{Coincidencia de patrones}
	% 			\paragraph{Constructores}
	% 			\paragraph{Resultados}
	% 			\paragraph{Autocompletado}
	% 			\paragraph{Búsqueda Avanzada}
	% 		\subsubsection{Implementación}
	% 			\paragraph{Investigación}
	% 			\paragraph{Estrategias}
	% 			\paragraph{Diseño}
	% \section{Técnicas de optimización para reducir tiempo de carga en sitios}
	% 	\subsection{Renderizado de contenido por navegadores}
	% 	\subsection{Optimización JS}
	% 		\paragraph{Optimización manual}
	% 		\paragraph{Minifiers}
	% 		\paragraph{Funciones inline}
	% 		\paragraph{DOM}
	% 		\paragraph{Delegación de eventos}
	% 		\paragraph{Ciclos}
	% 		\paragraph{Scope}
	% 		\paragraph{Async}
	% 	\subsection{Optimización CSS}
	% 		\paragraph{Propiedades abreviadas}
	% 		\paragraph{Agrupación de selectores}
	% 		\paragraph{Selectores anidados}
	% 		\paragraph{Minifiers}
	% 	\subsection{Optimización HTML}
	% 		\paragraph{Comentarios y espacios}
	% 		\paragraph{CSS y JS de línea}
	% 		\paragraph{Atributos Booleanos}
	% 	\subsection{Optimización de imágenes}
	% 		\paragraph{Formatos}
	% 		\paragraph{Optimización por zona}
	% 		\paragraph{Lazy Loadinf}
	% 		\paragraph{Renderizado de imágenes}
	% 		\paragraph{Data URIs}
	% 		\paragraph{Ancho y Alto}
	% 		\paragraph{Mapeo de imágenes}
	% 	\subsection{Otros tips}
	% 		\paragraph{CSS en head}
	% 	\subsection{Optimización PHP}
	% 		\paragraph{Optimización menor}
	% 		\paragraph{Optimización mayor}
	% 			\begin{enumerate}
	% 				\item Ciclos
	% 				\item Regex
	% 				\item Incluir archivos
	% 			\end{enumerate}
	% 	\subsection{Lighthouse}
\end{document}







































